% GENERATED FILE. Edit canonical lessons, then run npm run generate:books.
% canonical-source-hash: fnv1a64:9abecc46ba7778b5
% canonical-lessons: AR-C375-maa-dhalika, AR-C375-illa-anna, AR-C375-amma, AR-C375-kadhalika, AR-C375-ilawatan-ala

\chapter{Contrast and addition}
\label{ch:ar-contrast-and-addition}
\hlchaptermodality{375}

\begin{chapteropening}
\textbf{By the end of this chapter:} \emph{I can say even so, yet, as for, likewise and moreover.}

Complete the last lesson of chapter 375: I can say even so, yet, as for, likewise and moreover.
\end{chapteropening}

\section[\texorpdfstring{\ar{مع} \ar{ذلك} (maʿa dhālika)}{مع ذلك (maʿa dhālika)}]{\ar{مع} \ar{ذلك} (maʿa dhālika) — even so; however}

\label{lesson:AR-C375-maa-dhalika}



\begin{quote}
Before the new one: say the Arabic for given that, then the Arabic for in order to.
\end{quote}

\subsection*{You'll want to know: \ar{مع} \ar{ذلك}}
\textbf{\ar{مع} \ar{ذلك}} — \emph{maʿa dhālika} — \textquotedblleft{}even so; however\textquotedblright{}.

\begin{quote}
\ar{كان} \ar{متعبا،} \ar{مع} \ar{ذلك} \ar{أكمل} \ar{عمله}.
\end{quote}

\begin{quote}
\emph{kāna mutʿaban, maʿa dhālika akmala ʿamalah.}
\end{quote}

\begin{quote}
He was tired; even so, he finished his work.
\end{quote}

Often with \ar{و}: \ar{ومع} \ar{ذلك}, and yet.

\begin{cousinweb}
\textbf{\ar{مع} \ar{ذلك}} (\emph{maʿa dhālika}) joins two native Arabic words, \emph{maʿa}, \textquotedblleft{}with\textquotedblright{}, and \emph{dhālika}, \textquotedblleft{}that\textquotedblright{}: with all that, even so. The deeper origin of such particles is not settled.
\end{cousinweb}

\begin{grammarlens}[title={What you've built}]
One more word for giving reasons and opinions.
\end{grammarlens}

\subsection*{Your turn}
\begin{itemize}
\raggedright
  \item \emph{Say it:} \emph{maʿa dhālika}
  \item \emph{Say it:} \emph{maʿa dhālika}, once more
  \item \emph{Say it:} \emph{kāna mutʿaban, maʿa dhālika akmala ʿamalah.} — He was tired; even so, he finished his work.
  \item \emph{Recall:} say the Arabic for given that, then the Arabic for in order to, then say \emph{maʿa dhālika} again
\end{itemize}

\begin{tcolorbox}[breakable,colback=teal!4,colframe=teal!35!black,title={Before you move on}]
What does \ar{مع} \ar{ذلك} mean? (\textquotedblleft{}Even so; however\textquotedblright{}.) Say it once more.
\end{tcolorbox}
\section[\texorpdfstring{\ar{إلا} \ar{أن} (illā anna)}{إلا أن (illā anna)}]{\ar{إلا} \ar{أن} (illā anna) — yet; but}

\label{lesson:AR-C375-illa-anna}



\begin{quote}
Before the new one: say the Arabic for in order to, then the Arabic for even so.
\end{quote}

\subsection*{You'll want to know: \ar{إلا} \ar{أن}}
\textbf{\ar{إلا} \ar{أن}} — \emph{illā anna} — \textquotedblleft{}yet; but\textquotedblright{}.

\begin{quote}
\ar{الفندق} \ar{جميل،} \ar{إلا} \ar{أن} \ar{الغرف} \ar{صغيرة}.
\end{quote}

\begin{quote}
\emph{al-funduq jamīl, illā anna al-ghuraf ṣaghīra.}
\end{quote}

\begin{quote}
The hotel is beautiful, but the rooms are small.
\end{quote}

More formal than \ar{لكن}.

\begin{cousinweb}
\textbf{\ar{إلا} \ar{أن}} (\emph{illā anna}) joins the native Arabic particles \emph{illā}, \textquotedblleft{}except\textquotedblright{}, and \emph{anna}, \textquotedblleft{}that\textquotedblright{}: except that, and so yet, but. The deeper origin of such particles is not settled.
\end{cousinweb}

\begin{grammarlens}[title={What you've built}]
One more word for giving reasons and opinions.
\end{grammarlens}

\subsection*{Your turn}
\begin{itemize}
\raggedright
  \item \emph{Say it:} \emph{illā anna}
  \item \emph{Say it:} \emph{illā anna}, once more
  \item \emph{Say it:} \emph{al-funduq jamīl, illā anna al-ghuraf ṣaghīra.} — The hotel is beautiful, but the rooms are small.
  \item \emph{Recall:} say the Arabic for in order to, then the Arabic for even so, then say \emph{illā anna} again
\end{itemize}

\begin{tcolorbox}[breakable,colback=teal!4,colframe=teal!35!black,title={Before you move on}]
What does \ar{إلا} \ar{أن} mean? (\textquotedblleft{}Yet; but\textquotedblright{}.) Say it once more.
\end{tcolorbox}
\section[\texorpdfstring{\ar{أما} (ammā)}{أما (ammā)}]{\ar{أما} (ammā) — as for}

\label{lesson:AR-C375-amma}



\begin{quote}
Before the new one: say the Arabic for even so, then the Arabic for yet.
\end{quote}

\subsection*{You'll want to know: \ar{أما}}
\textbf{\ar{أما}} — \emph{ammā} — \textquotedblleft{}as for\textquotedblright{}.

\begin{quote}
\ar{أما} \ar{أخي} \ar{فيحب} \ar{الرياضة}.
\end{quote}

\begin{quote}
\emph{ammā akhī fa-yuḥibbu ar-riyāḍa.}
\end{quote}

\begin{quote}
As for my brother, he loves sport.
\end{quote}

The next part starts with \ar{فـ}: \ar{أما} … \ar{فـ} ….

\begin{cousinweb}
\textbf{\ar{أما}} (\emph{ammā}) is an old native Arabic particle that sets one topic apart from the others. Its origin is not settled.
\end{cousinweb}

\begin{grammarlens}[title={What you've built}]
One more word for giving reasons and opinions.
\end{grammarlens}

\subsection*{Your turn}
\begin{itemize}
\raggedright
  \item \emph{Say it:} \emph{ammā}
  \item \emph{Say it:} \emph{ammā}, once more
  \item \emph{Say it:} \emph{ammā akhī fa-yuḥibbu ar-riyāḍa.} — As for my brother, he loves sport.
  \item \emph{Recall:} say the Arabic for even so, then the Arabic for yet, then say \emph{ammā} again
\end{itemize}

\begin{tcolorbox}[breakable,colback=teal!4,colframe=teal!35!black,title={Before you move on}]
What does \ar{أما} mean? (\textquotedblleft{}As for\textquotedblright{}.) Say it once more.
\end{tcolorbox}
\section[\texorpdfstring{\ar{كذلك} (kadhālika)}{كذلك (kadhālika)}]{\ar{كذلك} (kadhālika) — likewise}

\label{lesson:AR-C375-kadhalika}



\begin{quote}
Before the new one: say the Arabic for yet, then the Arabic for as for.
\end{quote}

\subsection*{You'll want to know: \ar{كذلك}}
\textbf{\ar{كذلك}} — \emph{kadhālika} — \textquotedblleft{}likewise\textquotedblright{}.

\begin{quote}
\ar{أختي} \ar{تحب} \ar{الموسيقى،} \ar{وأمي} \ar{كذلك}.
\end{quote}

\begin{quote}
\emph{ukhtī tuḥibbu al-mūsīqā, wa-ummī kadhālika.}
\end{quote}

\begin{quote}
My sister loves music, and so does my mother.
\end{quote}

\begin{cousinweb}
\textbf{\ar{كذلك}} (\emph{kadhālika}) is built from the native Arabic \emph{ka-}, \textquotedblleft{}like\textquotedblright{}, and \emph{dhālika}, \textquotedblleft{}that\textquotedblright{}: like that, likewise. The deeper origin of such particles is not settled.
\end{cousinweb}

\begin{grammarlens}[title={What you've built}]
One more word for giving reasons and opinions.
\end{grammarlens}

\subsection*{Your turn}
\begin{itemize}
\raggedright
  \item \emph{Say it:} \emph{kadhālika}
  \item \emph{Say it:} \emph{kadhālika}, once more
  \item \emph{Say it:} \emph{ukhtī tuḥibbu al-mūsīqā, wa-ummī kadhālika.} — My sister loves music, and so does my mother.
  \item \emph{Recall:} say the Arabic for yet, then the Arabic for as for, then say \emph{kadhālika} again
\end{itemize}

\begin{tcolorbox}[breakable,colback=teal!4,colframe=teal!35!black,title={Before you move on}]
What does \ar{كذلك} mean? (\textquotedblleft{}Likewise\textquotedblright{}.) Say it once more.
\end{tcolorbox}
\section[\texorpdfstring{\ar{علاوة} \ar{على} (ʿilāwatan ʿalā)}{علاوة على (ʿilāwatan ʿalā)}]{\ar{علاوة} \ar{على} (ʿilāwatan ʿalā) — moreover; on top of}

\label{lesson:AR-C375-ilawatan-ala}



\begin{quote}
Before the new one: say the Arabic for as for, then the Arabic for likewise.
\end{quote}

\subsection*{You'll want to know: \ar{علاوة} \ar{على}}
\textbf{\ar{علاوة} \ar{على}} — \emph{ʿilāwatan ʿalā} — \textquotedblleft{}moreover; on top of\textquotedblright{}.

\begin{quote}
\ar{علاوة} \ar{على} \ar{الراتب،} \ar{يحصل} \ar{الموظف} \ar{على} \ar{سيارة}.
\end{quote}

\begin{quote}
\emph{ʿilāwatan ʿalā ar-rātib, yaḥṣulu al-muwaẓẓaf ʿalā sayyāra.}
\end{quote}

\begin{quote}
On top of the salary, the employee gets a car.
\end{quote}

Also \ar{علاوة} alone: a bonus on a salary.

\begin{cousinweb}
\textbf{\ar{علاوة} \ar{على}} (\emph{ʿilāwatan ʿalā}) is built on the native Arabic root \emph{ʿ-l-w}, \textquotedblleft{}to be high, to rise above\textquotedblright{} (Lane). \emph{ʿIlāwa} is what is laid on top of a load: something extra, and so moreover.
\end{cousinweb}

\begin{grammarlens}[title={What you've built}]
One more word for giving reasons and opinions.
\end{grammarlens}

\subsection*{Your turn}
\begin{itemize}
\raggedright
  \item \emph{Say it:} \emph{ʿilāwatan ʿalā}
  \item \emph{Say it:} \emph{ʿilāwatan ʿalā}, once more
  \item \emph{Say it:} \emph{ʿilāwatan ʿalā ar-rātib, yaḥṣulu al-muwaẓẓaf ʿalā sayyāra.} — On top of the salary, the employee gets a car.
  \item \emph{Recall:} say the Arabic for as for, then the Arabic for likewise, then say \emph{ʿilāwatan ʿalā} again
\end{itemize}

\begin{tcolorbox}[breakable,colback=teal!4,colframe=teal!35!black,title={Before you move on}]
What does \ar{علاوة} \ar{على} mean? (\textquotedblleft{}Moreover; on top of\textquotedblright{}.) Say it once more.
\end{tcolorbox}
